\IfFileExists{stacks-project.cls}{%
\documentclass{stacks-project}
}{%
\documentclass{amsart}
}

% For dealing with references we use the comment environment
\usepackage{verbatim}
\newenvironment{reference}{\comment}{\endcomment}
%\newenvironment{reference}{}{}
\newenvironment{slogan}{\comment}{\endcomment}
\newenvironment{history}{\comment}{\endcomment}

% For commutative diagrams we use Xy-pic
\usepackage[all]{xy}

% We use 2cell for 2-commutative diagrams.
\xyoption{2cell}
\UseAllTwocells

% We use multicol for the list of chapters between chapters
\usepackage{multicol}

% This is generally recommended for better output
\usepackage{lmodern}
\usepackage[T1]{fontenc}

% For cross-file-references
\usepackage{xr-hyper}

% Package for hypertext links:
\usepackage{hyperref}

% For any local file, say "hello.tex" you want to link to please
% use \externaldocument[hello-]{hello}
\externaldocument[introduction-]{introduction}
\externaldocument[conventions-]{conventions}
\externaldocument[sets-]{sets}
\externaldocument[categories-]{categories}
\externaldocument[topology-]{topology}
\externaldocument[sheaves-]{sheaves}
\externaldocument[sites-]{sites}
\externaldocument[stacks-]{stacks}
\externaldocument[fields-]{fields}
\externaldocument[algebra-]{algebra}
\externaldocument[brauer-]{brauer}
\externaldocument[homology-]{homology}
\externaldocument[derived-]{derived}
\externaldocument[simplicial-]{simplicial}
\externaldocument[more-algebra-]{more-algebra}
\externaldocument[smoothing-]{smoothing}
\externaldocument[modules-]{modules}
\externaldocument[sites-modules-]{sites-modules}
\externaldocument[injectives-]{injectives}
\externaldocument[cohomology-]{cohomology}
\externaldocument[sites-cohomology-]{sites-cohomology}
\externaldocument[dga-]{dga}
\externaldocument[dpa-]{dpa}
\externaldocument[sdga-]{sdga}
\externaldocument[hypercovering-]{hypercovering}
\externaldocument[schemes-]{schemes}
\externaldocument[constructions-]{constructions}
\externaldocument[properties-]{properties}
\externaldocument[morphisms-]{morphisms}
\externaldocument[coherent-]{coherent}
\externaldocument[divisors-]{divisors}
\externaldocument[limits-]{limits}
\externaldocument[varieties-]{varieties}
\externaldocument[topologies-]{topologies}
\externaldocument[descent-]{descent}
\externaldocument[perfect-]{perfect}
\externaldocument[more-morphisms-]{more-morphisms}
\externaldocument[flat-]{flat}
\externaldocument[groupoids-]{groupoids}
\externaldocument[more-groupoids-]{more-groupoids}
\externaldocument[etale-]{etale}
\externaldocument[chow-]{chow}
\externaldocument[intersection-]{intersection}
\externaldocument[pic-]{pic}
\externaldocument[weil-]{weil}
\externaldocument[adequate-]{adequate}
\externaldocument[dualizing-]{dualizing}
\externaldocument[duality-]{duality}
\externaldocument[discriminant-]{discriminant}
\externaldocument[derham-]{derham}
\externaldocument[local-cohomology-]{local-cohomology}
\externaldocument[algebraization-]{algebraization}
\externaldocument[curves-]{curves}
\externaldocument[resolve-]{resolve}
\externaldocument[models-]{models}
\externaldocument[functors-]{functors}
\externaldocument[equiv-]{equiv}
\externaldocument[pione-]{pione}
\externaldocument[etale-cohomology-]{etale-cohomology}
\externaldocument[proetale-]{proetale}
\externaldocument[relative-cycles-]{relative-cycles}
\externaldocument[more-etale-]{more-etale}
\externaldocument[trace-]{trace}
\externaldocument[crystalline-]{crystalline}
\externaldocument[spaces-]{spaces}
\externaldocument[spaces-properties-]{spaces-properties}
\externaldocument[spaces-morphisms-]{spaces-morphisms}
\externaldocument[decent-spaces-]{decent-spaces}
\externaldocument[spaces-cohomology-]{spaces-cohomology}
\externaldocument[spaces-limits-]{spaces-limits}
\externaldocument[spaces-divisors-]{spaces-divisors}
\externaldocument[spaces-over-fields-]{spaces-over-fields}
\externaldocument[spaces-topologies-]{spaces-topologies}
\externaldocument[spaces-descent-]{spaces-descent}
\externaldocument[spaces-perfect-]{spaces-perfect}
\externaldocument[spaces-more-morphisms-]{spaces-more-morphisms}
\externaldocument[spaces-flat-]{spaces-flat}
\externaldocument[spaces-groupoids-]{spaces-groupoids}
\externaldocument[spaces-more-groupoids-]{spaces-more-groupoids}
\externaldocument[bootstrap-]{bootstrap}
\externaldocument[spaces-pushouts-]{spaces-pushouts}
\externaldocument[spaces-chow-]{spaces-chow}
\externaldocument[groupoids-quotients-]{groupoids-quotients}
\externaldocument[spaces-more-cohomology-]{spaces-more-cohomology}
\externaldocument[spaces-simplicial-]{spaces-simplicial}
\externaldocument[spaces-duality-]{spaces-duality}
\externaldocument[formal-spaces-]{formal-spaces}
\externaldocument[restricted-]{restricted}
\externaldocument[spaces-resolve-]{spaces-resolve}
\externaldocument[formal-defos-]{formal-defos}
\externaldocument[defos-]{defos}
\externaldocument[cotangent-]{cotangent}
\externaldocument[examples-defos-]{examples-defos}
\externaldocument[algebraic-]{algebraic}
\externaldocument[examples-stacks-]{examples-stacks}
\externaldocument[stacks-sheaves-]{stacks-sheaves}
\externaldocument[criteria-]{criteria}
\externaldocument[artin-]{artin}
\externaldocument[quot-]{quot}
\externaldocument[stacks-properties-]{stacks-properties}
\externaldocument[stacks-morphisms-]{stacks-morphisms}
\externaldocument[stacks-limits-]{stacks-limits}
\externaldocument[stacks-cohomology-]{stacks-cohomology}
\externaldocument[stacks-perfect-]{stacks-perfect}
\externaldocument[stacks-introduction-]{stacks-introduction}
\externaldocument[stacks-more-morphisms-]{stacks-more-morphisms}
\externaldocument[stacks-geometry-]{stacks-geometry}
\externaldocument[moduli-]{moduli}
\externaldocument[moduli-curves-]{moduli-curves}
\externaldocument[examples-]{examples}
\externaldocument[exercises-]{exercises}
\externaldocument[guide-]{guide}
\externaldocument[desirables-]{desirables}
\externaldocument[coding-]{coding}
\externaldocument[obsolete-]{obsolete}
\externaldocument[fdl-]{fdl}
\externaldocument[index-]{index}

% Theorem environments.
%
\theoremstyle{plain}
\newtheorem{theorem}[subsection]{Theorem}
\newtheorem{proposition}[subsection]{Proposition}
\newtheorem{lemma}[subsection]{Lemma}

\theoremstyle{definition}
\newtheorem{definition}[subsection]{Definition}
\newtheorem{example}[subsection]{Example}
\newtheorem{exercise}[subsection]{Exercise}
\newtheorem{situation}[subsection]{Situation}

\theoremstyle{remark}
\newtheorem{remark}[subsection]{Remark}
\newtheorem{remarks}[subsection]{Remarks}

\numberwithin{equation}{subsection}

% Macros
%
\def\lim{\mathop{\mathrm{lim}}\nolimits}
\def\colim{\mathop{\mathrm{colim}}\nolimits}
\def\Spec{\mathop{\mathrm{Spec}}}
\def\Hom{\mathop{\mathrm{Hom}}\nolimits}
\def\Ext{\mathop{\mathrm{Ext}}\nolimits}
\def\SheafHom{\mathop{\mathcal{H}\!\mathit{om}}\nolimits}
\def\SheafExt{\mathop{\mathcal{E}\!\mathit{xt}}\nolimits}
\def\Sch{\mathit{Sch}}
\def\Mor{\mathop{\mathrm{Mor}}\nolimits}
\def\Ob{\mathop{\mathrm{Ob}}\nolimits}
\def\Sh{\mathop{\mathit{Sh}}\nolimits}
\def\NL{\mathop{N\!L}\nolimits}
\def\CH{\mathop{\mathrm{CH}}\nolimits}
\def\proetale{{pro\text{-}\acute{e}tale}}
\def\etale{{\acute{e}tale}}
\def\QCoh{\mathit{QCoh}}
\def\Ker{\mathop{\mathrm{Ker}}}
\def\Im{\mathop{\mathrm{Im}}}
\def\Coker{\mathop{\mathrm{Coker}}}
\def\Coim{\mathop{\mathrm{Coim}}}

% Boxtimes
%
\DeclareMathSymbol{\boxtimes}{\mathbin}{AMSa}{"02}

%
% Macros for moduli stacks/spaces
%
\def\QCohstack{\mathcal{QC}\!\mathit{oh}}
\def\Cohstack{\mathcal{C}\!\mathit{oh}}
\def\Spacesstack{\mathcal{S}\!\mathit{paces}}
\def\Quotfunctor{\mathrm{Quot}}
\def\Hilbfunctor{\mathrm{Hilb}}
\def\Curvesstack{\mathcal{C}\!\mathit{urves}}
\def\Polarizedstack{\mathcal{P}\!\mathit{olarized}}
\def\Complexesstack{\mathcal{C}\!\mathit{omplexes}}
% \Pic is the operator that assigns to X its picard group, usage \Pic(X)
% \Picardstack_{X/B} denotes the Picard stack of X over B
% \Picardfunctor_{X/B} denotes the Picard functor of X over B
\def\Pic{\mathop{\mathrm{Pic}}\nolimits}
\def\Picardstack{\mathcal{P}\!\mathit{ic}}
\def\Picardfunctor{\mathrm{Pic}}
\def\Deformationcategory{\mathcal{D}\!\mathit{ef}}
\usepackage{cleveref}
\crefname{originalreference}{original reference}{original references}
\Crefname{originalreference}{Original reference}{Original references}
\makeatletter
\newlabel{lemma-describe-sheaves-simplicial-site}{{09VM}{source}{}{sourcepointer-1}{}}
\newlabel{lemma-describe-sheaves-simplicial-site@cref}{{[originalreference][0][]09VM}{[1][0][]source}{}{}{}}
\makeatother
\crefname{originalreference}{original reference}{original references}
\Crefname{originalreference}{Original reference}{Original references}
\makeatletter
\newlabel{lemma-restriction-to-components}{{09W0}{source}{}{sourcepointer-2}{}}
\newlabel{lemma-restriction-to-components@cref}{{[originalreference][0][]09W0}{[1][0][]source}{}{}{}}
\makeatother
\crefname{originaldefinition}{original definition}{original definitions}
\Crefname{originaldefinition}{Original definition}{Original definitions}
\makeatletter
\newlabel{simplicial-definition-augmentation}{{018G}{source}{}{sourcepointer-3}{}}
\newlabel{simplicial-definition-augmentation@cref}{{[originaldefinition][0][]018G}{[1][0][]source}{}{}{}}
\makeatother
\title{Canonical morphism of topoi Sh(Y\_Zar) -> Sh(X) attached to an augmentation of simplicial spaces}
\author{The Stacks Project Authors}
\date{stacks-project:spaces-simplicial.tex:local-snapshot-20261012; GNU FDL 1.2 or later}
\begin{document}
\maketitle
\noindent\textit{Editorial note.} All delivered mathematical text is the author's own, copied from the pinned Stacks Project chapter file sources/stacks-spaces-simplicial/pinned-original-source.tex; no statement or proof step was rewritten, reordered or removed.\par\medskip
\noindent\textit{Editorial note.} The counted proof is the author's complete argument: the adjunction bijection is constructed in both directions (the component construction and the equalizer construction) and the author omits only the verification that they are mutually inverse ('We omit the proof that the constructions given above define mutually inverse bijections', source lines 578-579) - accepted per the precedent of card L46149 (restriction-to-components), where the same pattern (both constructions explicit, inverse check omitted) was documented; the closing remark that 'this morphism of topoi does not correspond to any obvious geometric functor between the sites defining the topoi' (source lines 584-586) is preserved verbatim as the author's observation.\par\medskip
\noindent\textit{Editorial note.} The label lemma-augmentation names the counted statement, and the proof environment immediately following it in the source is delivered verbatim as its proof.\par\medskip
\noindent\textit{Editorial note.} The counted statement is carried by the author's own lemma environment in the chapter source.\par\medskip
\noindent\textit{Allowed standard prerequisites (editorial).}\begin{itemize}
\item Augmentations of simplicial spaces. Simplicial, Definition \textbackslash{}label\{simplicial-definition-augmentation\}. Statement lines 525-527: Y is a simplicial space with augmentation a : Y -> X, with component morphisms of topological spaces a\_n : Y\_n -> X.
\item Sheaves on the simplicial site described componentwise. Same chapter, \textbackslash{}label\{lemma-describe-sheaves-simplicial-site\}. Proof lines 549-552: (1) is used as the definition of a\textasciicircum{}\{-1\}F via the componentwise description of sheaves on the simplicial site; (4) is used as the definition of a\_*.
\item The component restrictions g\_n : Sh(X) -> Sh(Y\_n). Same chapter, \textbackslash{}label\{lemma-restriction-to-components\} (card L46149). Statement lines 537-538 and proof lines 552-554: part (3), a g\_n = a\_n, holds because g\_n\textasciicircum{}\{-1\} a\textasciicircum{}\{-1\} = a\_n\textasciicircum{}\{-1\} by construction.
\end{itemize}\medskip
\ifdefined\hypertarget\hypertarget{sourcepointer-1}{}\fi
\noindent\textit{Source reference.} Original reference 09VM: sources/stacks-spaces-simplicial/pinned-original-source.tex, source lines 352--352. Prerequisite pointer invoked at proof line(s) 551; it names a standard fact of the same chapter that the counted argument uses, and every step of the counted proof built on it is retained verbatim in the delivered body.\par\medskip
\ifdefined\hypertarget\hypertarget{sourcepointer-2}{}\fi
\noindent\textit{Source reference.} Original reference 09W0: sources/stacks-spaces-simplicial/pinned-original-source.tex, source lines 456--456. Prerequisite pointer invoked at proof line(s) 538; it names a standard fact of the same chapter that the counted argument uses, and every step of the counted proof built on it is retained verbatim in the delivered body.\par\medskip
\ifdefined\hypertarget\hypertarget{sourcepointer-3}{}\fi
\noindent\textit{Source reference.} Original definition 018G: sources/stacks-chow/tags.txt, source lines 1615--1615. Prerequisite pointer invoked at proof line(s) 526; it names a standard fact of a sibling chapter that the counted argument uses, and every step of the counted proof built on it is retained verbatim in the delivered body.\par\medskip
\begin{lemma}
\label{lemma-augmentation}
Let $Y$ be a simplicial space and let $a : Y \to X$ be an augmentation
(Simplicial, Definition \ref{simplicial-definition-augmentation}).
Let $a_n : Y_n \to X$ be the corresponding morphisms of topological spaces.
There is a canonical morphism of topoi
$$
a : \Sh(Y_{Zar}) \to \Sh(X)
$$
with the following properties:
\begin{enumerate}
\item $a^{-1}\mathcal{F}$ is the sheaf restricting to $a_n^{-1}\mathcal{F}$
on $Y_n$,
\item $a_m \circ Y(\varphi) = a_n$ for all $\varphi : [m] \to [n]$,
\item $a \circ g_n = a_n$ as morphisms of topoi with
$g_n$ as in Lemma \ref{lemma-restriction-to-components},
\item $a_*\mathcal{G}$ for $\mathcal{G} \in \Sh(Y_{Zar})$
is the equalizer of the two maps
$a_{0, *}\mathcal{G}_0 \to a_{1, *}\mathcal{G}_1$.
\end{enumerate}
\end{lemma}

\begin{proof}
Part (2) holds for augmentations of simplicial objects in any category.
Thus $Y(\varphi)^{-1} a_m^{-1} \mathcal{F} = a_n^{-1}\mathcal{F}$
which defines an $Y(\varphi)$-map from $a_m^{-1}\mathcal{F}$
to $a_n^{-1}\mathcal{F}$.
Thus we can use (1) as the definition of $a^{-1}\mathcal{F}$ (using
Lemma \ref{lemma-describe-sheaves-simplicial-site}) and
(4) as the definition of $a_*$. If this defines a morphism of topoi
then part (3) follows because we'll have $g_n^{-1} \circ a^{-1} = a_n^{-1}$
by construction. To check $a$ is a morphism of topoi we have to show
that $a^{-1}$ is left adjoint to $a_*$ and we have to show that
$a^{-1}$ is exact. The last fact is immediate from the exactness of
the functors $a_n^{-1}$.

\medskip\noindent
Let $\mathcal{F}$ be an object of $\Sh(X)$ and let $\mathcal{G}$
be an object of $\Sh(Y_{Zar})$. Given
$\beta : a^{-1}\mathcal{F} \to \mathcal{G}$ we can look at the
components $\beta_n : a_n^{-1}\mathcal{F} \to \mathcal{G}_n$.
These maps are adjoint to maps
$\beta_n : \mathcal{F} \to a_{n, *}\mathcal{G}_n$.
Compatibility with the simplicial structure shows that
$\beta_0$ maps into $a_*\mathcal{G}$.
Conversely, suppose given a map $\alpha : \mathcal{F} \to a_*\mathcal{G}$.
For any $n$ choose a $\varphi : [0] \to [n]$. Then we can look at
the composition
$$
\mathcal{F} \xrightarrow{\alpha} a_*\mathcal{G}
\to a_{0, *}\mathcal{G}_0 \xrightarrow{\mathcal{G}(\varphi)}
a_{n, *}\mathcal{G}_n
$$
These are adjoint to maps $a_n^{-1}\mathcal{F} \to \mathcal{G}_n$
which define a morphism of sheaves $a^{-1}\mathcal{F} \to \mathcal{G}$.
We omit the proof that the constructions given above define
mutually inverse bijections
$$
\Mor_{\Sh(Y_{Zar})}(a^{-1}\mathcal{F}, \mathcal{G}) =
\Mor_{\Sh(X)}(\mathcal{F}, a_*\mathcal{G})
$$
This finishes the proof. An interesting observation is here that
this morphism of topoi does not correspond to any obvious geometric
functor between the sites defining the topoi.
\end{proof}

\end{document}
